\section{Introduction}

Autoregressive decoding in large language models (LLMs) is inherently
sequential: each generated token requires a full forward pass of the target
model. This constraint directly limits single-stream throughput in interactive
inference settings. Speculative decoding mitigates this bottleneck by adding a
smaller draft model that proposes a length-$k$ token block, followed by one
target-side verification step; modified rejection sampling preserves the target
distribution~\cite{Leviathan2023,Chen2023specsampling}.

Although speculative decoding is now well studied, deployment guidance remains
uncertain for consumer-GPU settings. Reported gains in the literature are often
obtained on data-center hardware and are sensitive to model pairing, proposal
length, and sampling regime. For practitioners operating on a single RTX-class
GPU, the key questions are practical: which draft size is worthwhile, which
$k$ is optimal, and whether deterministic versus stochastic decoding changes
the speed--quality frontier.

This gap is an engineering pain point: speedups demonstrated on H100-class
servers do not transfer directly to mobile RTX deployments, where memory
bandwidth, launch overhead, and power-constrained duty cycle can dominate
decode-time behavior.

This direction is increasingly important because real-world LLM usage is moving
toward consumer-class deployment constraints: local assistants, student and
developer laptops, and privacy-sensitive workloads where cloud offload is
undesirable or expensive. In this setting, efficiency is not only a cost
optimization problem; it is an accessibility and feasibility requirement.
Inference methods that are robust on consumer hardware can widen practical LLM
adoption beyond data-center environments.

\subsection{Research questions}
\label{sec:rqs}

We address three research questions:
\begin{enumerate}\itemsep0pt
  \item[\textbf{RQ1.}] How much end-to-end speedup does vanilla
        speculative decoding deliver against autoregressive decoding on a
        single RTX~4090 with a 3\,B target, and how close is this to
        a practical speedup ceiling under consumer-GPU constraints?
  \item[\textbf{RQ2.}] How do draft model size and draft length jointly
        determine token acceptance and net latency, and what does this reveal
        about verification-loop overhead and effective GPU duty cycle?
  \item[\textbf{RQ3.}] How sensitive is the speed--quality trade-off to
        the sampling regime (deterministic vs.\ stochastic) and to task
        entropy?
\end{enumerate}

\subsection{Contributions}
\label{sec:contributions}

This paper makes the following contributions:
\begin{enumerate}\itemsep0pt
  \item A reproducible consumer-GPU protocol with frozen sample manifests,
        matched prompts, and two decoding regimes across GSM8K, MMLU subset,
        and CNN/DailyMail.
  \item A complete 12-configuration speculative grid (two draft sizes,
        $k\in\{4,8,16\}$, deterministic/stochastic) with regime-matched
        baselines and per-configuration reporting of latency, throughput,
        acceptance, and quality deltas.
  \item Empirical evidence that, in the current RTX4090 CUDA-graph run family,
        speedup decreases with larger $k$ (highest at $k{=}4$), with best
        observed performance at 0.5B, $k{=}4$, deterministic ($S{=}3.0674$),
        and a deterministic advantage in 5 of 6 matched draft--$k$ pairs.
  \item A statistical reporting protocol for camera-ready evaluation: paired
        bootstrap confidence intervals, one-sided significance tests for
        $S{>}1$, and corrected pairwise comparisons for winner selection.
  \item An empirical ACSD extension run (adaptive $k\in\{3,4,5\}$ with
        0.5B$\rightarrow$1.5B rescue drafting) that approaches best fixed-policy
        speedup while adding explicit tail-control guardrails.
\end{enumerate}

The rest of the paper is organised as follows. Section~\ref{sec:related}
reviews related work. Section~\ref{sec:protocol} defines the experimental
protocol; Section~\ref{sec:metrics} defines metrics. Section~\ref{sec:impl}
describes implementation and reproducibility controls. Section~\ref{sec:results}
reports the main empirical findings. Section~\ref{sec:acsd} outlines
our ACSD empirical extension beyond fixed-policy drafting. Section~\ref{sec:discussion}
discusses implications, limitations, and conclusion.
